\section{Conclusion and Further Research}
In this study, I used UKHLS data on chronic illnesses and daily ailments to identify three health types: the ``resilient health type'' (Type~I) is considered to be the common type accounting for 92\% of the clustering sample, the ``early onset frailty type'' (Type~II) suffers from higher initial frailty but a steadier decline after age 60, and the ``accelerated decline type'' (Type~III) begins with decent health similar to the Resilient Type but experiences a rapid decline in health starting the middle age and retirement age.

The identification of these distinct health trajectories within an economic sample has significant implications for retirement policy and economic inequality. My analysis of employment rates, educational attainment, and earnings across these health types reveals substantial disparities. The resilient health type consistently demonstrates higher employment rates and earnings, as well as a higher proportion of degree holders. In contrast, the early onset frailty and accelerated decline types show lower labour market participation and earnings, suggesting a potential cycle of health-related economic disadvantage.

These findings underscore the need for nuanced policy approaches that account for diverse health trajectories. Traditional retirement policies based on a uniform retirement age may inadequately address the needs of individuals in the Early Onset Frailty or Accelerated Decline groups, who may require earlier interventions or more flexible work arrangements. Moreover, the stark differences in economic outcomes across health types highlight the potential for health-driven economic inequality, emphasizing the importance of targeted healthcare and social support initiatives.

Further research should delve deeper into the underlying biological factors that determine an individual's health type. The UKHLS dataset offers opportunities using biomarkers, genetic data, and proteomic information. Incorporating these biological indicators as propensity score variables could provide insights beyond the impact of age-related factors, potentially uncovering ``economic markers'' of an individual's health trajectory. Such markers could serve as early warning signs, allowing for preemptive interventions in both healthcare and economic planning.

Additionally, longitudinal studies tracking individuals' transitions between health types over time could offer valuable insights into the malleability of these trajectories and the effectiveness of various interventions. Exploring the interplay between occupation types, lifestyle factors, and health trajectories could further illuminate the complex relationship between work, health, and economic outcomes.

Finally, I observed the effect of health type as exogenous variables with large explanatory power in regard to predicting one’s frailty. Although as \citet{denardi2024wp} mentioned, the formation of health types are certainly the result of earlier build-ups before adulthood \citep{borella_nber}. The combined effects should be manifested well beyond the analysis through significant disparities in retirement savings, lifetime earnings, and marriage patterns across health types. This underscores the importance of considering health trajectories in economic models and policy formulation, as they not only influence individual well-being but also have far-reaching implications for societal productivity, healthcare costs, and economic inequality. Understanding these factors would be tremendously beneficial to policymakers and individuals alike.
